Ein Lastwagen (\mLO) mit \vLO\ und ein Auto (\mAO) mit \vAO\ stossen frontal zusammen.

\begin{center}
\begin{tikzpicture}[>=Stealth]
 \draw[thick] (-0.3,0) -- (9.3,0);
 % truck, facing right
 \draw[fill=gray!25] (0,0.25) rectangle (2.8,1.5);
 \draw[fill=gray!40] (2.85,0.25) -- (2.85,1.2) -- (3.4,1.2) -- (3.75,0.7) -- (3.75,0.25) -- cycle;
 \foreach \x in {0.5,1.2,3.2} \draw[fill=white] (\x,0.22) circle (0.22);
 \node at (1.4,0.9) {Lastwagen};
 \draw[->, very thick] (1.2,1.8) -- (2.6,1.8) node[right] {\vLO};
 % car, facing left
 \draw[fill=gray!15] (5.4,0.25) -- (5.4,0.6) -- (5.9,0.65) -- (6.3,1.0) -- (7.3,1.0) -- (7.7,0.65) -- (7.8,0.25) -- cycle;
 \foreach \x in {5.9,7.3} \draw[fill=white] (\x,0.2) circle (0.2);
 \node[font=\small] at (6.7,0.5) {Auto};
 \draw[->, very thick] (7.2,1.35) -- (5.8,1.35) node[left] {\vAO};
\end{tikzpicture}
\end{center}

\begin{abcliste}
\abc
Vergleiche während des Zusammenstosses die Kraft des Lastwagens auf das Auto mit der Kraft des Autos auf den Lastwagen (Betrag und Richtung).

\abc
Tim sagt: „Der Lastwagen drückt stärker, weil er schwerer ist.“ Nimm Stellung.

\abc
Während des Zusammenstosses wirkt zwischen den Fahrzeugen eine mittlere Kraft von \FKO. Berechne die (mittleren) Beschleunigungen der beiden Fahrzeuge. Andere Kräfte sind daneben vernachlässigbar.

\abc
Wessen Geschwindigkeit ändert sich beim Zusammenstoss stärker, und wievielmal?
\end{abcliste}
